\documentclass[conference]{IEEEtran}
\IEEEoverridecommandlockouts

\usepackage{cite}
\usepackage{amsmath,amssymb,amsfonts}
\usepackage{graphicx}
\usepackage{textcomp}
\usepackage{xcolor}
\usepackage{booktabs}
\usepackage{multirow}
\usepackage{array}
\usepackage{tikz}
\usepackage{pgfplots}
\pgfplotsset{compat=1.18}
\usepackage[hidelinks]{hyperref}
\usepackage{xr-hyper}
\externaldocument[][nocite]{main}

\newcommand{\status}[1]{\textsc{#1}}
\renewcommand{\thesection}{S\Roman{section}}
\renewcommand{\thetable}{S\Roman{table}}
\renewcommand{\thefigure}{S\arabic{figure}}
\renewcommand{\theequation}{S\arabic{equation}}

\begin{document}

\title{Supplementary Material for\\ ``Agentic Pilot: An Evidence-Driven Local Autonomous AI Agent Framework\\for Privacy-Preserving Intelligent Task Automation''}
\author{}
\maketitle

This supplement holds the material that does not fit the paper's page limit: the full simulation methodology, results and parameters (all \textsc{synthetic} unless marked), the end-to-end harness (implemented, not executed), the design of the unimplemented computer backend, and the threat table. References to sections, tables and figures without the prefix ``S'' point to the paper.

\section{Gate Ablation}
\label{sec:abl-full}

\begin{table}[t]
\centering
\caption{Ablation of the gate on all 92 constructed states: false completions (FC) and false rejections (FR) with one conjunct of $\Phi$, or the bot-check halt $\beta$, removed. $p$: exact McNemar test of correct decisions against the full gate on the design states (paired by state); descriptive only, since the conjuncts were written on these states. On the held-out sets no removal changes more than three decisions. Source: \texttt{verification\_gate\_ablation.csv}.}
\label{tab:ablation}
\footnotesize
\setlength{\tabcolsep}{3pt}
\begin{tabular}{@{}lccccccc@{}}
\toprule
 & \multicolumn{3}{c}{Design (23/15)} & \multicolumn{2}{c}{Held-out 1 (5/7)} & \multicolumn{2}{c}{Held-out 2 (22/20)} \\
\cmidrule(lr){2-4}\cmidrule(lr){5-6}\cmidrule(l){7-8}
Removed & FC & FR & $p$ & FC & FR & FC & FR \\
\midrule
none (full $\Phi$) & 3 & 1 & -- & 3 & 1 & 13 & 1 \\
$E$ error, sign-in & 8 & 1 & .06 & 4 & 1 & 14 & 1 \\
$H$ search home & 7 & 1 & .13 & 3 & 1 & 15 & 1 \\
$Q$ query match & 4 & 1 & 1.0 & 3 & 1 & 13 & 1 \\
$T$ typed text & 6 & 1 & .25 & 3 & 1 & 14 & 1 \\
$X$ answer & 6 & 0 & .63 & 3 & 1 & 13 & 1 \\
$R$ pending steps & 4 & 1 & 1.0 & 3 & 1 & 15 & 1 \\
$\beta$ bot check & 5 & 1 & .50 & 4 & 0 & 15 & 0 \\
\bottomrule
\end{tabular}
\end{table}


Table~\ref{tab:ablation} removes one conjunct at a time. On the design states every removal raises false completions, most for the error check $E$ (3$\to$8) and the search-home check $H$ (3$\to$7), but with 23 unsatisfied states no single removal changes more than five decisions, and with at most five discordant states the smallest attainable exact McNemar $p$ is 0.063, so the test could not reach $p<0.05$; this is not evidence that a conjunct is unnecessary. On the held-out sets each removal changes at most three decisions, and removing $Q$ or $X$ changes none: these conjuncts rarely fire on phrasings they were not written for. Only two removals trade false completions for false rejections: the bot-check halt $\beta$ (held-out~2: 13$\to$15 FC, 1$\to$0 FR) and, on the design states, the answer check $X$ (3$\to$6 FC, 1$\to$0 FR). Evidence types thus contribute on the states that trigger them, and the gate's coverage, not the strength of any single check, limits it.


\section{Simulation Methodology}
\label{sec:sim-full}

This section gives the full methodology summarized in Section~\ref{sec:sim} of the paper.

\paragraph{Why simulation} The end-to-end protocol needs local model weights, which the environment used for this paper could not download. A simulation is appropriate for the remaining questions because the quantities they depend on separate into component behaviour that we can measure without a model (detection and false-rejection rates of $\Phi$, latencies, the engine's classification, retrieval ranks) and planner behaviour that we cannot; propagating the first through the control flow of the loop, while sweeping the second, shows which unmeasured parameters decide H1--H5. We built such a seeded simulation (\texttt{benchmarks/simulation}). The simulation does not execute the agent. Its outputs follow from measured and assumed inputs. Every number in this paper belongs to one of six classes: \textsc{actual} (genuinely executed or measured on the implementation), \textsc{synthetic} (produced by the simulation), \textsc{assumed} (an input value chosen by us, not measured), \textsc{derived} (computed from existing results, e.g.\ a rate estimated from actual counts), \textsc{literature} (taken from cited work), or \textsc{proposed} (behaviour of a component that is not implemented, or an evaluation not yet done). Tables and figures state the class of their values in the caption or row label; in the text, numbers in Sections~\ref{sec:results} and~\ref{sec:safety} are \textsc{actual} unless stated, and numbers in Section~\ref{sec:sim-results} are \textsc{synthetic} unless marked otherwise. An evidence ledger lists every simulated value in the main tables with its source and class.

\paragraph{Measured inputs} Without a model we measured, on the implementation (2026-10-23, 4-core Intel Xeon at 2.8\,GHz, same software versions): the decisions of $\Phi$ and of a basic success detector (error check $E$ only) on the 92 gate states; the latency of navigation, clicks, form filling, DOM extraction, screenshots and $\Phi$ on the self-hosted sites (30 repetitions each); the failure class and strategy sequence that the real recovery engine assigns to six failure messages; and the ranks at which the real Chroma memory store returns relevant records for ten queries over a 22-record synthetic memory corpus. Detection rates used by the simulation are derived from the two held-out gate sets only (for detection, the 22 unsatisfied states left after the bot-check category is excluded, 18 from held-out~2 and 4 from held-out~1; the 24 satisfied states of the same categories for false rejection), with the estimator $(k+0.5)/(n+1)$, which gives a small nonzero rate (0.07--0.13) where none was observed, as for basic detection in five categories and $\Phi$ on extraction, and turns the 0 of 24 false rejections into 0.02. Simulated categories map to gate categories as follows: navigation, search, extraction and multi-step to themselves, data entry to text entry, and form submission to text entry and multi-step pooled. The detection rate of a missing confirmation is proxied by the multi-step rate (an assumed proxy; 3 states), and the category mapping above is likewise assumed.

\paragraph{Assumptions and generation} What needs a model---how often the planner claims completion of an unfinished task, how often actions fail and are recovered, how much memory helps---is assumed and listed with its class in Table~\ref{tab:sim-params} (Section~\ref{app:simparams}). A \emph{claim} in the simulation stands for any point at which the code evaluates $\Phi$: a \texttt{complete} action, any successful action on the last plan step (for a request that is not decomposed, every successful action), or the typing of requested exact text; $\pi$ is the probability that such a point is reached while the task is unfinished. A scenario has a task category (navigation, search, form submission, data entry, extraction, multi-step) and an outcome class (successful, incomplete, incorrect action, missing evidence, contradictory evidence, stale evidence, ambiguous) that sets the probability that the goal is actually met. Each scenario draws its random numbers once, and every configuration reuses them, so comparisons are paired by scenario. The true outcome is evaluated per configuration because memory and continuation after a rejection change the final state. For v1, parameters and code were committed (\texttt{3a0aa8e}) before the held-out set (seed 202, 100 scenarios) was generated; the reported results come from the post-hoc v2 described below; a development set (seed 101) was used only to debug the code. After the freeze, one function was changed: the latency model (Study~F) did not give all configurations the same random draws, so it was corrected to common random numbers in \texttt{41a6097}; Studies A--E and H were not affected and gave identical outputs before and after. Later, \texttt{episode()} gained three output fields (claim, true state at the claim, first decision) for the decision-level analysis; outputs were otherwise byte-identical.

\paragraph{Corrected model (v2, post hoc)} A methodology audit of the frozen model (v1) found that it contradicted the implementation in four places: its false-rejection rates came from bot-check pages, a category the scenarios do not contain; a rejected claim followed by a successful continuation was accepted without a second check, although the verification node checks every claim; the basic detector did not fail closed on an unreadable page, although it shares the error check $E$; and recovery was applied to every action failure, although action errors are first retried three times inline and a missing confirmation reaches the recovery engine only as a budget failure after 15 model calls. The results reported here come from a corrected model, v2 (\texttt{simulate\_v2.py}), that fixes these four points, draws derived rates from their Beta posteriors in the Monte Carlo replications, adds re-planning to the latency model and sweeps the continuation rate. v2 reuses v1's generator, seeds, scenario set and assumed values. It was written after the v1 results were known and is therefore \emph{not pre-registered}; the corrections moved the results for $\Phi$ in its favour (task success relative to no verification from 37 against 43, $p=0.07$, to 44 against 46; decision accuracy from 45/71 against 43/71 to 52/74 against 46/74), mainly because v1's false-rejection rate (0.089) came from bot-check pages; v1's outputs are archived (\texttt{out/heldout\_v1}) and every number that changed is listed in the statistical correction report (\texttt{artifacts/results/final\_audit/}). Monte Carlo replications regenerate the scenario set 1{,}000 times and draw the derived rates from their posteriors, to show the variability the generator and the rate uncertainty induce; the assumed values and the recovery probabilities taken from Study~E are held fixed, so the intervals omit their uncertainty.

\paragraph{Metrics} With TP, FP, FN, TN the confusion counts of reported completion against the true outcome over $N$ scenarios: $\mathrm{TSR}=\mathrm{TP}/N$; FCR and FRR as in Eq.~\eqref{eq:fcr}, so that precision $=1-\mathrm{FCR}$ and recall $=1-\mathrm{FRR}$; F1 and accuracy $(\mathrm{TP}+\mathrm{TN})/N$. Recovery success is recovered over recoverable failures; retrieval quality is precision and recall at $k=3$ and the mean reciprocal rank (MRR) of the first relevant record. Episode time is decomposed as
\begin{equation}
T = T_{\mathrm{inf}} + T_{\mathrm{browser}} + T_{\mathrm{evidence}} + T_{\mathrm{verify}} + T_{\mathrm{recovery}},
\label{eq:time}
\end{equation}
with browser, evidence and $\Phi$ times resampled from the measurements and inference time assumed. Two levels are kept apart. \emph{Policy-level} outcomes describe how an episode ends (task success, episode-level false completion); \emph{decision-level} outcomes describe whether the verifier's first decision on a completion claim matches the true state at that claim. Because a rejection can change how an episode ends, the final ground truth of a scenario may differ between configurations, so only the decision level compares verifiers against one shared truth. The decision-level outcome was defined after the v1 results were known, when an audit showed that the earlier accuracy comparison mixed the two levels. Paired configurations are compared by the exact McNemar test on per-scenario binary indicators of one level, and the variability of paired differences across regenerated scenario sets by Monte Carlo percentiles. Every $p$-value and interval computed from synthetic data is conditional on the simulation assumptions; it describes the synthetic sample, not the deployed agent.

\paragraph{Studies} Study~A generates the scenarios and computes the completion metrics, B compares four completion rules, C ablates components, D evaluates memory, E recovery, F overhead, G privacy (analytically, Section~\ref{sec:safety}) and H sensitivity.

\paragraph{Validation by recomputation} For v1, a separate AI agent of the same authoring pipeline, given only the per-scenario output files, recomputed, with separately written code, the confusion counts, rates and paired tests of Studies B and C and the memory and recovery statistics from the per-scenario files, and found no discrepancy (\texttt{simulation\_validation\_report.md}). Monte Carlo, sensitivity, latency, decision-level and recovery paired-test values were not independently recomputed. This check is a recomputation of the analysis code's arithmetic from the same synthetic files; it is neither independent human annotation of any label nor an independent replication of an experiment. For v2, a final AI review agent of the same pipeline recomputed with separately written code the policy- and decision-level confusion counts, the paired tests and the recovery and memory values from the per-scenario files, reran the simulation and the table generator byte for byte, and found no discrepancy; Monte Carlo, sensitivity and latency values were checked only through that rerun; the recomputation scripts are in \texttt{artifacts/results/final\_audit/scripts}.


\section{Simulation Results}
\label{sec:sim-results-full}

This section gives the full results summarized in Section~\ref{sec:sim-results} of the paper; Fig.~\ref{fig:sim-sens} is in the paper.

\noindent\fbox{\parbox{0.96\columnwidth}{\small All values in this section are \textsc{synthetic} unless marked \textsc{actual} or stated as \textsc{assumed} inputs; rows marked $^{\ddagger}$ are \textsc{proposed} components simulated with assumed rates. The synthetic values come from the corrected seeded simulation (v2, post hoc) of Section~\ref{sec:sim}, not from running the agent. They show the consequences of the stated assumptions and the measured component behaviour; they do not establish how the deployed agent performs. Every $p$-value and interval here is conditional on those assumptions.}}

\subsection{Baselines and Verification (Studies A and B; H1, H2)}

\begin{table*}[t]
\centering
\caption{\textsc{Simulation-based (synthetic)} Study~B: completion decisions of four configurations on the same 100 held-out synthetic scenarios. Confusion counts and rates with Wilson 95\% intervals. $^{\ddagger}$Not implemented (proposed).}
\label{tab:sim-baseline}
\footnotesize
\setlength{\tabcolsep}{5pt}
\begin{tabular}{@{}lccccccccc@{}}
\toprule
Configuration & TP & FP & FN & TN & TSR & FCR & FRR & F1 & Acc. \\
\midrule
No verification & 43 & 28 & 3 & 26 & 43/100 [0.34,0.53] & 28/71 [0.29,0.51] & 3/46 [0.02,0.18] & .735 & .690 \\
Basic success detection & 42 & 24 & 4 & 30 & 42/100 [0.33,0.52] & 24/66 [0.26,0.48] & 4/46 [0.03,0.20] & .750 & .720 \\
Evidence-gated ($\Phi$) & 37 & 19 & 10 & 34 & 37/100 [0.28,0.47] & 19/56 [0.23,0.47] & 10/47 [0.12,0.35] & .718 & .710 \\
Enhanced validation$^{\ddagger}$ & 39 & 8 & 12 & 41 & 39/100 [0.30,0.49] & 8/47 [0.09,0.30] & 12/51 [0.14,0.37] & .796 & .800 \\
\bottomrule
\end{tabular}
\end{table*}


\begin{table}[t]
\centering
\caption{\textsc{Simulation-based (synthetic)} Paired comparisons with the evidence-gated configuration, exact McNemar tests on discordant pairs $b/c$ (first count: only the other configuration; second: only $\Phi$). FC: episode ends in a false completion (100 scenarios). Decision: first decision correct (71 claims). Success: task achieved and reported (100 scenarios). MC: mean FCR and 2.5--97.5th percentiles over 1000 regenerated scenario sets. Synthetic $p$-values and intervals are conditional on the simulation assumptions; they are not empirical tests of the agent.}
\label{tab:sim-tests}
\scriptsize
\setlength{\tabcolsep}{1.8pt}
\begin{tabular}{@{}lcccccc@{}}
\toprule
Configuration & FC $b/c$ & $p$ & Decision $b/c$ & $p$ & Success $b/c$; $p$ & MC FCR \\
\midrule
No verification & 9/0 & 0.004 & 7/9 & 0.804 & 7/1; 0.070 & 0.36 [0.25,0.48] \\
Basic success detection & 5/0 & 0.062 & 6/5 & 1.000 & 6/1; 0.125 & 0.33 [0.22,0.44] \\
Evidence-gated ($\Phi$) & ref. & ref. & ref. & ref. & ref. & 0.25 [0.15,0.37] \\
Enhanced validation$^{\ddagger}$ & 0/11 & 0.001 & 11/2 & 0.022 & 4/2; 0.688 & 0.14 [0.05,0.24] \\
\bottomrule
\end{tabular}
\end{table}


Table~\ref{tab:sim-baseline} compares four completion rules on the same 100 scenarios of the reported synthetic set (seed 202) at two levels that must not be mixed. At the \emph{decision level}, the planner claims completion in 74 scenarios, 46 of them truly complete at that moment; the four verifiers face the same 74 claims and the same truth, so their decisions are paired. At the \emph{policy level}, an episode's final state may differ between configurations for the same scenario, because a rejected claim returns control to planning and the episode may then finish (or, after a false rejection, never be reported); task success and episode-level false completions are paired outcomes of different policies, not of one ground truth.

\emph{H1 (false completions).} Episodes end in a false completion 19 times with $\Phi$ against 28 without verification (FCR 0.30 against 0.38; discordant pairs 9 and 0; Table~\ref{tab:sim-tests}); basic success detection ends in 24. In the model the direction is fixed by construction (in the code, a continuation after a false rejection could later be accepted on a changed page, which the model omits): $\Phi$ can only withhold a claim that would otherwise be accepted, so no scenario can end in a false completion with $\Phi$ alone, one discordant cell is always 0, and only the size of the reduction is informative. Across 1{,}000 regenerated scenario sets, with the derived rates drawn from their posteriors, the reduction averages 0.11 false completions per scenario (2.5--97.5th percentiles 0.05--0.18); because one discordant cell is structurally 0, a McNemar $p$ would only restate the 9 discordant scenarios, so we report this interval instead. Under the stated assumptions the simulation is consistent with H1; it is not empirical confirmation.

\emph{H2 (decision reliability).} On the 74 claims, $\Phi$ decides correctly in 52 (accuracy 0.70, Wilson 0.59--0.79) and accepting every claim in 46 (0.62, 0.51--0.72): $\Phi$ is correct alone in 9 claims, all rejections of premature claims, and wrong alone in 3, all false rejections (two of them on scenarios with missing evidence, which $\Phi$ rejects at an \textsc{assumed} rate of 0.8; the measured false-rejection rate is 0 of 24) (exact McNemar $p=0.15$). The difference has the hypothesized direction but is not significant in this set, and across regenerated sets $\Phi$ is better at $p<0.05$ in only 42\% (mean accuracy 0.74 against 0.64), so H2 is not supported for $\Phi$ at this sample size (74 claims). The proposed enhanced validation reaches 61 of 74 (0.82; against $\Phi$ 11 and 2, $p=0.022$), a difference that follows from its assumed detection rates. At the policy level, the simulated $\Phi$ configuration ends with 44 of 100 synthetic scenarios achieved and reported, against 46 without verification (3 only without verification, 1 only with $\Phi$; $p=0.63$): in the simulation, false rejections cost 3 successes and continuation after a correct rejection recovers 1.

\subsection{Ablation (Study C)}

\begin{table}[t]
\centering
\caption{\textsc{Simulation-based (synthetic)} Study~C: one component removed at a time, on the same 100 scenarios (policy level). Each row is a simulated configuration, not a run of the implementation; the rows for $\Phi$, memory and recovery correspond to switches that exist in the code (\texttt{verification\_mode}, \texttt{enable\_memory}/\texttt{enable\_rag}, \texttt{enable\_recovery}). $^{\S}$Simulated removal of all stored evidence (screenshots and records); the implemented switch \texttt{enable\_evidence} removes only the per-action execution records. Decisions are unchanged because $\Phi$ reads the live page, not the stored evidence (code behaviour). Evid.: share of episodes with audit evidence stored. MC: means over 1000 regenerated sets.}
\label{tab:sim-ablation}
\scriptsize
\setlength{\tabcolsep}{2pt}
\begin{tabular}{@{}lccccccc@{}}
\toprule
Configuration & TSR & FCR & FRR & F1 & Evid. & MC TSR & MC FCR \\
\midrule
Full & 37/100 [0.28,0.47] & .339 & .213 & .718 & 0.89 & 0.44 & 0.25 \\
-- stored evidence$^{\S}$ & 37/100 [0.28,0.47] & .339 & .213 & .718 & 0.00 & 0.44 & 0.25 \\
-- $\Phi$ & 43/100 [0.34,0.53] & .394 & .065 & .735 & 0.89 & 0.45 & 0.36 \\
-- memory, RAG & 38/100 [0.29,0.48] & .345 & .191 & .724 & 0.89 & 0.44 & 0.26 \\
-- recovery & 30/100 [0.22,0.40] & .348 & .250 & .698 & 0.90 & 0.38 & 0.25 \\
\bottomrule
\end{tabular}
\end{table}


Table~\ref{tab:sim-ablation} removes one component at a time. Removing $\Phi$ reproduces the no-verification row. In the simulation, disabling the recovery engine, with the retries before it still in place, lowers synthetic TSR from 0.44 to 0.41 (3 successes lost, none gained; $p=0.25$) and, across replications, from 0.47 to 0.45; the size of this effect is set by the recovery probabilities of Study~E, which rest on assumed per-attempt success rates. Removing memory changes one outcome in each direction and nothing on average. Stored evidence is not ablated: $\Phi$ reads the live page, so stored evidence serves the audit trail and does not change decisions.

\subsection{Memory and Retrieval (Study D; H3)}

\begin{table}[t]
\centering
\caption{Study~D, memory and retrieval. Upper block ACTUAL (real store, synthetic records); lower block \textsc{Simulation-based (synthetic)}. P(misleading): probability that the retrieved context misleads the planner; for the store it is set to the measured share of irrelevant top-1 records (DERIVED; that this share is the probability of misleading the planner is ASSUMED); for the filter it is an assumed input, so any effect of the filter follows from that assumption. With 44 repeated workflows the study cannot detect an effect of the assumed size (gain 0.08, loss 0.10). $^{\ddagger}$Not implemented.}
\label{tab:sim-memory}
\scriptsize
\setlength{\tabcolsep}{3pt}
\begin{tabular}{@{}lcccc@{}}
\toprule
\multicolumn{5}{@{}l}{\emph{ACTUAL retrieval, real store, 22 synthetic records; 9 queries with a relevant record}} \\
Precision@3 / Recall@3 / MRR & 0.52 & 1.00 & 0.93 &  \\
Top-1 irrelevant or failed (all 10 queries) & 3/10 &  &  &  \\
\midrule
\multicolumn{5}{@{}l}{\emph{SYNTHETIC repeated workflows (n=44)}} \\
Configuration & TSR & FCR & P(misleading) &  \\
Memory disabled & 19/44 [0.30,0.58] & .387 & -- &  \\
Memory enabled & 19/44 [0.30,0.58] & .367 & 0.30 (der.) &  \\
Relevance-filtered$^{\ddagger}$ & 20/44 [0.32,0.60] & .355 & 0.10 (assumed) &  \\
\bottomrule
\end{tabular}
\end{table}


On the synthetic memory corpus the real store returned a relevant record in the top three for each of the 9 queries that had one (recall@3 1.00, MRR 0.93 over those 9; \textsc{actual}). It has no relevance threshold, however: a query with no relevant memory still returned three records, and for 3 of 10 queries the top record was irrelevant or a \emph{failed} earlier attempt (Table~\ref{tab:sim-memory}). Relevance is labelled by task group and is loose (a search for one place counts as relevant to a search for another), so recall is optimistic, and precision@3 is bounded by the number of relevant records per query; the measured 0.52 is the largest value attainable on this corpus. With a misleading rate set to the measured share of irrelevant top-1 records (3/10; that this share equals the probability of misleading the planner is assumed), the simulated repeated workflows succeed 19 times of 44 with and without memory, so H3 is not supported. With 44 workflows and an assumed gain of 0.08 the study could not have detected an effect of that size, so it is uninformative about H3 rather than evidence against it. The proposed relevance filter (20 of 44) changes the outcome only through its assumed misleading rate.

\subsection{Recovery (Study E; H4)}

\begin{table}[t]
\centering
\caption{\textsc{Simulation-based (synthetic)} Study~E: 120 failure scenarios, 20 per mode; 27 drawn as permanent (25\% expected), none of them a missing confirmation. Entry paths and strategy sequences follow the code: action errors are retried three times inline before the engine is called; navigation failures and requests for help go to the engine directly; a missing confirmation raises no error and is handled only by rejection and continuation. Per-attempt success rates are ASSUMED (PROPOSED for the guided variant), so differences between rows follow from them. Detected: failure noticed; FC: undetected missing confirmations reported complete; Unnec.: attempts on permanent failures; Att.: mean attempts. $^{\ddagger}$Not implemented.}
\label{tab:sim-recovery}
\scriptsize
\setlength{\tabcolsep}{2pt}
\begin{tabular}{@{}lccccc@{}}
\toprule
Recovery & Detected & Recovered & FC & Unnec. & Att. \\
\midrule
Inline retries only & 111/120 [0.86,0.96] & 29/93 [0.23,0.41] & 9 & 51 & 1.28 \\
Retry + engine (impl.) & 111/120 [0.86,0.96] & 67/93 [0.62,0.80] & 9 & 162 & 3.12 \\
Evidence-guided$^{\ddagger}$ & 118/120 [0.94,1.00] & 87/93 [0.87,0.97] & 2 & 162 & 3.17 \\
\midrule
\multicolumn{6}{@{}l}{\emph{Recovered / recoverable by mode (inline only; retry + engine; guided)}} \\
\multicolumn{6}{@{}l}{nav. 0; 14; 16 / 16, time-out 9; 12; 13 / 14, missing el. 9; 13; 14 / 15} \\
\multicolumn{6}{@{}l}{invalid sub. 9; 13; 14 / 14, unexp. state 0; 13; 13 / 14, missing conf. 2; 2; 17 / 20} \\
\bottomrule
\end{tabular}
\end{table}


The real recovery engine classifies a navigation error containing \texttt{net::ERR\_} as transient, treats an invalid submission as a verification failure and assigns it re-planning, and does not recognize a missing confirmation, which it classifies as unknown (\textsc{actual}). In the code, action errors, including a planner's request for help, are first retried three times inline and only then passed to the engine; a failed initial navigation is retried three times by the browser executor before it reaches the engine; and a missing confirmation raises no error at all. Simulated along these paths with assumed per-attempt success rates (Table~\ref{tab:sim-recovery}), the retries before the engine recover 51 of 93 synthetic recoverable failures, and adding the engine raises this to 68 of 93. Missing confirmations are recovered only by continuation after a rejection (2 of 20 with or without the engine), and 9 of them end as false completions. The proposed evidence-guided recovery, which checks an explicit confirmation postcondition, recovers 87 and leaves 2 false completions (19 scenarios recovered only by it, none only by the engine); this is a consequence of its assumed success rates, not evidence for H4. In the simulation the engine spends 174 attempts on permanent failures, the guided variant 162.

\subsection{Verification Overhead (Study F; H5)}

\begin{table}[t]
\centering
\caption{Study~F, verification overhead. Upper block ACTUAL; lower block \textsc{Simulation-based (synthetic)}, built from the ACTUAL component samples and an ASSUMED inference time; the two blocks are not combined into one claim. Inference dominates every synthetic episode (84\% of time under $\Phi$); evidence screenshots take 2.5\%. $^{\ddagger}$Not implemented.}
\label{tab:sim-overhead}
\scriptsize
\setlength{\tabcolsep}{3pt}
\begin{tabular}{@{}lccc@{}}
\toprule
\multicolumn{4}{@{}l}{\emph{ACTUAL component latency (ms), local sites, 30 repetitions}} \\
Operation & median & p95 &  \\
navigate & 17.1 & 24.9 &  \\
click + load & 65.2 & 106.2 &  \\
fill & 12.5 & 17.4 &  \\
DOM extraction & 76.4 & 105.6 &  \\
screenshot (evidence) & 61.3 & 93.6 &  \\
$\Phi$ on a site page & 5.2 & 6.9 &  \\
$\Phi$ on gate states (83) & 26.8 & 30.3 &  \\
\midrule
\multicolumn{4}{@{}l}{\emph{SYNTHETIC episode time (s), inference median 3.0 s assumed}} \\
Configuration & median & p95 & verif.\ share \\
No verification & 17.0 & 34.1 & 0.0\% \\
Evidence-gated ($\Phi$) & 17.0 & 34.2 & 0.2\% \\
LLM self-verification & 21.7 & 38.7 & 19.6\% \\
$\Phi$ + model check$^{\ddagger}$ & 21.7 & 38.7 & 19.8\% \\
\bottomrule
\end{tabular}
\end{table}


$\Phi$ takes a median of 5.2\,ms on an ordinary page of the self-hosted sites and 26.8\,ms on the constructed gate states; a screenshot 61\,ms and DOM extraction 76\,ms (\textsc{actual}, Table~\ref{tab:sim-overhead}). In the simulation, with an assumed median of 3\,s per model call on a CPU, inference takes 80\% of an episode and the checks of $\Phi$ 0.1\%. The mean episode is nevertheless 2.6\% longer with $\Phi$ (20.2 against 19.7\,s), because each rejected claim costs one further planning step (planner call, DOM extraction, click, three screenshots; assumed), and $\Phi$ times are resampled from the slower gate-state measurements; a model self-check adds 14\% (22.4\,s; a cost-only figure, since the simulated self-check accepts every claim). At 1.5\,s or 6\,s per call the median episode under $\Phi$ takes 10.0\,s or 36.8\,s. H5 is consistent with the simulation: the rule itself is cheap, and its cost is in false rejections and the re-planning that rejections cause.

\subsection{Robustness and Sensitivity (Study H)}

\begin{table}[t]
\centering
\caption{\textsc{Simulation-based (synthetic)} Study~H: robustness of the evidence-gated configuration (means over 300 regenerated sets per row, drawn from a random stream separate from the 1{,}000-set Monte Carlo of Table~\ref{tab:sim-tests}, so the unperturbed rows differ slightly from it). Without verification FCR stays at 0.36 at every continuation rate.}
\label{tab:sim-sens}
\scriptsize
\setlength{\tabcolsep}{3pt}
\begin{tabular}{@{}lccc@{}}
\toprule
Varied factor ($\Phi$) & FCR & FRR & TSR \\
\midrule
$\Phi$ detection $\times0.5$ & 0.30 & 0.14 & 0.46 \\
$\Phi$ detection $\times1.0$ & 0.25 & 0.14 & 0.47 \\
$\Phi$ detection $\times1.5$ & 0.19 & 0.13 & 0.49 \\
missing evidence 0\% & 0.23 & 0.09 & 0.52 \\
missing evidence 8\% & 0.25 & 0.14 & 0.47 \\
missing evidence 20\% & 0.29 & 0.22 & 0.39 \\
action failure rate 0.1 & 0.25 & 0.14 & 0.50 \\
action failure rate 0.25 & 0.25 & 0.14 & 0.47 \\
action failure rate 0.4 & 0.25 & 0.14 & 0.44 \\
finish after rejection 0.1 & 0.26 & 0.14 & 0.45 \\
finish after rejection 0.35 & 0.25 & 0.14 & 0.47 \\
finish after rejection 0.6 & 0.24 & 0.13 & 0.50 \\
\bottomrule
\end{tabular}
\end{table}


The advantage of verification depends on how often the planner claims completion prematurely, which is unmeasured (Fig.~\ref{fig:sim-sens}): at $\pi=0.3$ the FCR of $\Phi$ is 0.14 against 0.21 without verification, at $\pi=0.95$ 0.32 against 0.45, while FRR stays near 0.14 against 0.08. Halving the derived detection rates of $\Phi$ raises its FCR from 0.25 to 0.30, about halfway to no verification (0.36); missing evidence raises FRR (0.09 to 0.22) because the simulated $\Phi$ rejects a claim on an unreadable page with assumed probability 0.8; more action failures lower TSR; and a higher rate of finishing after a rejection raises the TSR of $\Phi$ from 0.45 to 0.50 (Table~\ref{tab:sim-sens}).

\section{Simulation Parameters and Provenance}
\label{app:simparams}

Table~\ref{tab:sim-params} lists the parameters of the simulation with their class; modelling choices that are not single values (latency structure, attempt caps) are listed at the end of the table. The assumed premature-claim rate (0.65) and finish-after-rejection rate (0.35) are swept in Study~H to show how far the results depend on them, and the rates that decide the gate's behaviour are derived from the held-out gate states, not from the design states on which the gate was tuned. The per-scenario data (\texttt{scenarios.csv}, \texttt{recovery\_scenarios.csv}), the results, the evidence ledger with the source and class of every simulated number, and the v1 validation report are in \texttt{benchmarks/simulation/out/heldout\_v2/} (v1: \texttt{out/heldout\_v1/}, identical to \texttt{out/heldout/}, which the v1 decision analysis reads) and \texttt{artifacts/results/}. The v2 outputs and the simulation tables are regenerated by \texttt{python -m benchmarks.simulation.simulate\_v2} and \texttt{make\_tables}, and the whole chain, including the gate statistics, the egress summary and the paper, by \texttt{make reproduce-paper} with the package versions of \texttt{requirements-paper.lock}; the v1 recomputation script is \texttt{independent\_validator.py} (an AI-written arithmetic check, not an independent replication).

\begin{table}[t]
\centering
\caption{Simulation parameters of the corrected model v2 and their provenance (\texttt{benchmarks/simulation/params.json}, \texttt{simulate\_v2.py}). DERIVED gate rates use $(k+0.5)/(n+1)$ on the held-out gate states (the misleading-retrieval rate is the raw share $3/10$); Monte Carlo sets draw them from Beta$(k+0.5, n-k+0.5)$.}
\label{tab:sim-params}
\scriptsize
\setlength{\tabcolsep}{2pt}
\begin{tabular}{@{}p{3.6cm}p{3.0cm}p{1.3cm}@{}}
\toprule
Parameter & Value & Class \\
\midrule
Scenarios per set; dev / held-out seed; MC sets & 100; 101 / 202; 1000 & ASSUMED (design choice) \\
Task categories (nav., search, form, data entry, extraction, multi-step) & 0.15 / 0.15 / 0.20 / 0.15 / 0.20 / 0.15 & ASSUMED \\
Outcome classes (successful, incomplete, incorrect, missing, contradictory, stale, ambiguous) & 0.40 / 0.15 / 0.12 / 0.08 / 0.08 / 0.07 / 0.10 & ASSUMED \\
P(task truly complete $\mid$ class), same order & 1.0 / 0.0 / 0.0 / 0.5 / 0.5 / 0.3 / 0.5 & ASSUMED \\
Planner claims completion: if complete / if not ($\pi$) & 0.92 / 0.65 & ASSUMED \\
$\Phi$ withhold rate on unsatisfied states, by task category & navi. 0.36, sear. 0.50, form. 0.50, data. 0.38, extr. 0.07, mult. 0.63 & DERIVED (mapping ASSUMED) \\
Basic detection withhold rate, same & navi. 0.36, sear. 0.10, form. 0.07, data. 0.13, extr. 0.07, mult. 0.13 & DERIVED (mapping ASSUMED) \\
False rejection: $\Phi$ / basic & 0/24 $\to$ 0.02 / 0.02 (satisfied states of the mapped categories; bot-check pages excluded) & DERIVED \\
Ambiguous-class detection factor & 0.5 & ASSUMED \\
Verifier rejects when evidence is missing ($\Phi$ and basic) & 0.8 & ASSUMED \\
Enhanced: extra catch / extra false rej. / contradiction rej. & 0.5 / 0.05 / 0.7 & PROPOSED \\
Finish after rejection (continuation is verified again) & 0.35 & ASSUMED \\
Action failure; recovered with / without engine & 0.25; 0.66 / 0.49 & ASSUMED; SYNTHETIC (Study E) \\
Repeated workflows; memory gain / loss & 0.4; 0.08 / 0.1 & ASSUMED \\
Misleading retrieval: store / filtered & 0.30 (raw $k/n$) / 0.1 & DERIVED (mapping ASSUMED) / PROPOSED \\
Recovery: permanent share; attempts & 0.25; 3 inline + up to 3 engine; 3 executor + up to 6 engine (navigation) & ASSUMED; code \\
Recovery entry paths; strategy sequences & code (Table~\ref{tab:sim-recovery}) & code inspection; ACTUAL \\
Recovery success per attempt & 0.10--0.45 (re-plan), 0.05--0.35 (vision); guided 0.35--0.55 & ASSUMED / PROPOSED \\
Missing-confirmation detection: retry / guided & 0.63 / 0.9 & DERIVED via ASSUMED proxy / PROPOSED \\
Steps per episode by category (same order as above) & 2 / 3 / 5 / 3 / 3 / 6 & ASSUMED \\
Inference time per call (median, log-sd) & 3.0 s, 0.5 & ASSUMED \\
Browser, DOM, screenshot latency; $\Phi$ latency & resampled from site measurements; from the gate-state measurements & ACTUAL \\
Category mapping (simulated to gate categories) & Section~\ref{sec:sim} & ASSUMED \\
Executor and inline retries & scored at the re-plan success rate & ASSUMED \\
Missing confirmation after rejection & one continuation attempt at the finish-after-rejection rate & ASSUMED \\
Episode time structure & 1 planner call, DOM extraction and action per step; 3 screenshots per step, 2 at completion; 1 extra step per rejection; recovery attempts from Study~E & ASSUMED (code-based) \\
Self-check baseline & one model call per claim; never rejects (cost only) & ASSUMED \\
Monte Carlo; sensitivity grid & 1000 sets, assumed values and recovery probabilities fixed; 300 sets per point & design \\
\bottomrule
\end{tabular}
\end{table}


\section{End-to-End Harness (Implemented, Not Executed)}
\label{sec:method-e2e-full}

\paragraph{Environment and tasks} Live websites change and block automated traffic, and WebArena-style site images were not available offline, so we built five self-hosted sites served on the loopback interface with fictional content that a model cannot answer from memory: a search engine, an encyclopedia, a shop with search and cart, a town-office contact form and newsletter, and a notes application; a further path redirects to a sign-in wall. The 40 tasks (10 development, 30 test) cover search, navigation, text entry, extraction, form submission, multi-step tasks and 5 unsatisfiable requests, at four difficulty tiers by the number of actions a solution needs (1, 2--3, 4--6, more). The notes site appears only in the test split. Tasks needing an account, payment, a CAPTCHA, the live web, or an outcome that neither server state nor answer text can verify were excluded by design. Because the search-home check $H$ recognizes only public search engines, it never fires on these sites; its ablation in the harness is uninformative.

\paragraph{Oracle} $\Phi^{\star}$ is computed from signals that neither the agent nor $\Phi$ uses: server-side records of the sites (cart contents, form submissions, saved notes, field values autosaved by page scripts), the final URL, and the ground-truth facts in the site data matched against the agent's answer; unsatisfiable tasks are never achieved. It calls no model and shares no code with the verifier. We measured its error rate on 114 scripted trajectories driven directly through Playwright without the agent: a correct solution for each task, an idle trajectory that only opens the start page, and a near miss for 34 tasks (wrong value, wrong item, extra submission, partial query).

\paragraph{Protocol} Table~\ref{tab:e2e} (Section~\ref{app:tasks}) lists the configurations: the full system, two baselines---(a) no completion check and (b) self-verification by the planner model---and one-factor ablations of each conjunct, memory, recovery and each recovery strategy. Each configuration runs in its own process with its own database and memory store, at temperature 0 with a fixed seed, three trials per task, and the same budget of 15 planner calls (after which the recovery engine may grant its remaining attempts). Per episode the harness records the oracle verdict, the agent's status, browser actions, model calls and tokens, wall-clock time, each verification decision with its latency, recovery strategies, retrieved memories, Ollama version and model digests, and the host of every page request Playwright reports. The analysis reports success, FCR and FRR with Wilson intervals, per-category breakdowns, exact McNemar tests and task-level bootstrap intervals~\cite{efron1979bootstrap} for differences to the full system, and the overhead of $\Phi$ relative to baseline (a). A published agent framework as a third baseline would need its own model and prompts and is outside the scope of this paper.


\subsection{Task Set}
\label{app:tasks}

The 40 tasks (\texttt{benchmarks/e2e/tasks.json}) are distributed as follows (development/test): extraction 2/8, form submission 2/7, multi-step 1/5, navigation 2/2, search 1/2, text entry 1/2, unsatisfiable 1/4; tiers 1--4 contain 3/6, 4/11, 3/12 and 0/1 tasks. The test split uses all five sites; the notes site appears only there. The harness (\texttt{run\_e2e.py}) takes the split, the number of trials and a list of the configurations of Table~\ref{tab:e2e}; the analysis script (\texttt{analyze\_e2e.py}) writes result tables with Wilson intervals, McNemar tests and bootstrap intervals.

\begin{table}[t]
\centering
\caption{End-to-end evaluation design (implemented in \texttt{benchmarks/e2e}, not yet run with a model). Each configuration changes one factor of the full system; all use the same tasks, model, sampling and budget.}
\label{tab:e2e}
\footnotesize
\setlength{\tabcolsep}{2.2pt}
\begin{tabular}{@{}p{2.95cm}p{2.85cm}p{2.6cm}@{}}
\toprule
Configuration & Change & Question \\
\midrule
\texttt{full} & -- & task success, FCR, FRR, cost \\
\texttt{no\_verification} & completion accepted wherever $\Phi$ would be evaluated & baseline (a); gain and overhead of $\Phi$ \\
\texttt{self\_verification} & planner model judges its own completion & baseline (b) \\
\texttt{gate\_without\_}$c$ & conjunct $c$ removed & per-evidence ablation \\
\texttt{no\_memory} & memory and RAG off & contribution of memory \\
\texttt{no\_evidence} & no per-action execution records & traceability \\
\texttt{no\_recovery} & engine off; inline retries remain & contribution of recovery \\
\texttt{recovery\_}$s$\texttt{\_only} & strategy $s$ pinned & per-strategy recovery \\
\texttt{single\_model} & routing off & routing \\
\texttt{vision\_only\_\allowbreak grounding} & no element list & grounding \\
\bottomrule
\end{tabular}
\end{table}




\section{Computer Backend (Design, Not Implemented)}
\label{app:computer}

The repository contains a desktop executor that wraps PyAutoGUI and psutil, a text-file helper and a file-existence predicate. None is connected to the control loop: no node, action type or API route calls them, and the file helper does not confine paths. The rest of this appendix is a design.

\paragraph{Observation and grounding} The backend would observe the foreground window through its accessibility tree (Windows UI Automation, macOS Accessibility, AT-SPI) into a manifest shaped like the DOM manifest, falling back to a vision model on a window screenshot. File and process actions are grounded in typed state (paths, hashes, process identity).

\paragraph{Postconditions and permission tiers} Each mutating action carries exact postconditions (Table~\ref{tab:post}) captured with before/after values. Because many local actions cannot be undone, a postcondition checked afterwards is no protection on its own, so an \texttt{action\_gate} node between planning and execution would compute a tier from the concrete action: T0 (read-only) runs freely; T1 (reversible, with a recorded inverse) runs automatically only if the user enabled it; T2 (destructive) and T3 (network-capable programs, credentials) always need approval bound to a hash of the exact action. Paths outside allow-listed roots, inline interpreters and privilege elevation would be refused. A design in which untrusted content cannot change the control flow~\cite{debenedetti2025camel} is a stronger alternative.

\begin{table}[t]
\centering
\caption{Completion postconditions by environment and their status in the code. Impl.: implemented and used by the agent; Unit: implemented and unit-tested, not called by the agent; Prop.: proposed design. $h_0$: hash before the action; $(\mathit{pid},t_c)$: process ID and creation time.}
\label{tab:post}
\footnotesize
\setlength{\tabcolsep}{2.2pt}
\begin{tabular}{@{}p{1.25cm}p{5.05cm}l@{}}
\toprule
Env. & Example postcondition & Status \\
\midrule
Browser & no error or sign-in page; results page for the quoted query; requested text present; answer mentions the goal; no pending step & Impl. \\
Browser & bot-check page detected $\Rightarrow$ \textsc{Blocked} & Impl. \\
File & path exists and size $\geq n$ & Unit \\
File & created / overwritten: $\mathrm{sha256}=\mathrm{sha256}(\text{content})$; moved: source absent and $\mathrm{sha256}(\text{dst})=h_0(\text{src})$; trashed: path absent and trash entry present & Prop. \\
Process & launched: $(\mathit{pid},t_c)$ alive; terminated: $(\mathit{pid},t_c)$ gone within a timeout & Prop. \\
Window & foreground title matches; window with title exists & Prop. \\
Desktop UI & accessibility value of a control equals the input; control state contains \textit{checked} & Prop. \\
Read-only & workspace manifest (paths, sizes, mtimes) unchanged & Prop. \\
\bottomrule
\end{tabular}
\end{table}



\section{Threat Surfaces and Security Limitations}
\label{sec:sec-full}

\begin{table}[t]
\centering
\caption{Privacy and security assessment (Study~G, analytical). Evidence: E egress test, C code inspection, T unit test; -- none.}
\label{tab:threats}
\scriptsize
\setlength{\tabcolsep}{2pt}
\begin{tabular}{@{}p{1.55cm}p{3.05cm}p{0.6cm}p{2.95cm}@{}}
\toprule
Surface & Control in the code & Evid. & Residual risk \\
\midrule
Model inference & loopback endpoint enforced; override must be explicit & E, T & user override sends prompts and screenshots off the machine \\
Memory, task data & SQLite and Chroma on local disk; telemetry off & C, T & not encrypted; readable by any local process \\
Websites & none: the planner may follow links to any site & C & data the agent types reaches that site \\
Browser traffic & page requests logged in the harness & E & Chromium background traffic not observed \\
Screenshots & stored in the local evidence folder & C & not redacted; may show account data \\
Retention & no deletion policy & C & evidence grows without bound \\
High-risk actions & approval before posting, sending, purchasing, deleting & C, T & model-assigned label; once per task \\
Injected text & 17-pattern sanitizer on element text & C, T & denylist; titles and memory unfiltered \\
\bottomrule
\end{tabular}
\end{table}


\paragraph{What it does not provide} The sanitizer is a fixed denylist, which systematic evaluations find insufficient~\cite{liu2024formalizing}; titles, extracted text and memory are not filtered. The action label and risk level that trigger approval come from a small model and are decided once per task, so an injected instruction that steers the planner to submit a form mid-task is not gated. The API on 127.0.0.1 has no authentication and checks neither the Host header nor the Origin of WebSocket connections, so a web page open on the same machine, including one the agent visits, may be able to reach it (DNS rebinding, cross-site WebSocket access) and read task events, records and evidence screenshots, subject to the browser's local-network protections; typed text is stored in clear, and screenshots are not redacted. We make no security claim. The browser hides common automation signals; we consider this acceptable only for a personal assistant acting at human pace on its user's requests, and deployments should disable it where a site's terms prohibit automated access.

\section{Implementation Audit}
\label{sec:status-full}

\begin{table}[t]
\centering
\caption{Implementation audit, traced through the code. Class: IE implemented and experimentally evaluated (component level, \textsc{actual}); IS implemented, studied only by simulation; IN implemented, not evaluated; PI partially implemented; PR proposed. Evidence: C component study, O oracle validation, E egress test with a stand-in model server, T unit test, M measured input to the simulation, S simulation.}
\label{tab:status}
\scriptsize
\setlength{\tabcolsep}{2.4pt}
\begin{tabular}{@{}p{4.3cm}cp{2.6cm}@{}}
\toprule
Component & Class & Evidence \\
\midrule
Completion predicate $\Phi$, bot-check halt & IE & C (92 states), M, S \\
Browser control (Playwright), DOM extraction & IE & C, M (latency) \\
Loopback-only model endpoint; Chroma telemetry off & IE & E (destinations, 70 stand-in episodes), T (refusal, telemetry flag) \\
End-to-end oracle $\Phi^{\star}$ & IE & O (114 trajectories) \\
End-to-end harness (run with a language model) & IN & E (stand-in only) \\
Recovery engine: classification, strategy order & IE & C (6 failure modes), S \\
Memory store retrieval (Chroma, SQLite) & IE & C (10 queries, synthetic corpus), S \\
Per-action evidence (screenshots, records) & IS & M (latency), S \\
LLM planner in the full loop & IN & tests with mocked model \\
Vision fallback; task-level human approval & IN & tests \\
Recovery: alternative selector, replan & PI & labels only \\
Self-verification baseline mode & IN & harness only \\
Enhanced validation; evidence-guided recovery; memory relevance filter & PR & S (assumed behaviour) \\
Desktop backend, action gate, permission tiers & PR & -- \\
\bottomrule
\end{tabular}
\end{table}



\bibliographystyle{IEEEtran}
\bibliography{references}

\end{document}
